\documentclass[11pt]{article}

\usepackage[T1]{fontenc}
\usepackage[utf8]{inputenc}
\usepackage{amsmath,amssymb,amsthm}
\usepackage[margin=1in]{geometry}
\usepackage[colorlinks=true,linkcolor=blue,citecolor=blue,urlcolor=blue]{hyperref}

\newtheorem{theorem}{Theorem}
\newtheorem{lemma}[theorem]{Lemma}

\title{Prime Summands of Surgeries on L--Space Knots}
\date{}

\begin{document}

\maketitle

\begin{abstract}
We prove that every closed 3--manifold obtained by Dehn surgery on an
L--space knot in \(S^3\) has at most two prime summands.  The key input is
Greene's theorem that the cabling conjecture holds for knots admitting
L--space surgeries.
\end{abstract}

\section*{Conventions}

All manifolds are compact, connected, and orientable.  For a knot
\(K \subset S^3\) and a slope \(r \in \mathbb Q \cup \{\infty\}\), write
\(S^3_r(K)\) for the result of \(r\)-surgery on \(K\).  We write slopes using
the meridian--longitude convention: \(a/b\) denotes the class
\(a\mu_K+b\lambda_K\), with \(\gcd(a,b)=1\).

By a prime summand we mean a non-\(S^3\) factor in the Kneser--Milnor prime
decomposition.  Thus \(S^3\)-summands are ignored, as usual.

An L--space knot means a knot in \(S^3\) admitting a nontrivial surgery to an
L--space.  This convention is invariant under mirroring.

\section*{Standard inputs}

We use the following two classical facts.

\begin{lemma}[Gordon--Luecke]\label{lem:nonintegral}
Let \(J \subset S^3\) be a knot.  If \(r \in \mathbb Q\) is nonintegral, then
\(S^3_r(J)\) is irreducible, except possibly when \(J\) is the unknot, in which
case \(S^3_r(J)\) is a lens space and is again irreducible unless it is
\(S^3\).
\end{lemma}

\begin{proof}
For nontrivial \(J\), this is precisely the theorem of Gordon--Luecke that a
reducible surgery on a knot in \(S^3\) must have integral slope
\cite[Theorem~1]{GL}.  For the unknot, nonzero surgeries are lens spaces; lens
spaces are irreducible except for \(L(1,q)\cong S^3\).
\end{proof}

\begin{lemma}\label{lem:lens-prime}
If \(p>1\), then every lens space \(L(p,q)\) is prime.
\end{lemma}

\begin{proof}
Lens spaces are irreducible; see, for example, Hempel
\cite[Chapter~2]{Hempel}.  Since \(L(p,q)\not\cong S^3\) for \(p>1\), an
irreducible lens space is prime.
\end{proof}

The deep theorem we need is Greene's proof of the cabling conjecture for
L--space knots.

\begin{theorem}[Greene: cabling conjecture for L--space knots]
\label{thm:greene}
Let \(K \subset S^3\) be a nontrivial knot admitting an L--space surgery.  If
some Dehn surgery on \(K\) is reducible, then \(K\) is a cable knot and the
reducing slope is the cabling slope.  Equivalently, for some knot
\(J \subset S^3\) and coprime integers \(p,q\) with \(p>1\),
\[
K = J_{p,q}
\]
is the \((p,q)\)-cable of \(J\), and the reducible surgery slope is \(pq\).
\end{theorem}

\begin{proof}[Reference]
This is Greene's theorem that the cabling conjecture holds for knots admitting
L--space surgeries; see Greene \cite[Main Theorem and its cabling-conjecture
corollary]{Greene}.  Torus knots are included by allowing the companion \(J\)
to be the unknot.
\end{proof}

We also recall the standard surgery formula for cables.

\begin{lemma}[Cable surgery formula]\label{lem:cable}
Let \(J \subset S^3\) be a knot and let \(J_{p,q}\) be its \((p,q)\)-cable,
where \(p>1\) and \(\gcd(p,q)=1\).  With the usual convention that \(p\) is the
winding number, the cabling slope is \(pq\), and
\[
S^3_{pq}(J_{p,q}) \cong L(p,q) \# S^3_{q/p}(J).
\]
\end{lemma}

\begin{proof}
This is the classical cable-surgery calculation; see Gordon
\cite{Gor83}.  Briefly, let \(V=\nu(J)\).  The cable \(J_{p,q}\) lies on
\(\partial V\) as a curve representing \(p\lambda_J+q\mu_J\).  The cabling
annulus has boundary slope \(pq\) on \(\partial \nu(J_{p,q})\).  After
\(pq\)-surgery this annulus caps off to a separating \(2\)-sphere.  One side is
the result of attaching a 2--handle to the solid torus \(V\) along a
\((p,q)\)-curve, hence is \(L(p,q)\) with a 3--ball removed.  The other side is
the exterior of \(J\) filled so that the meridian of the attached solid torus
is \(q\mu_J+p\lambda_J\), namely \(q/p\)-surgery on \(J\).  Thus the filled
manifold is
\[
L(p,q)\# S^3_{q/p}(J).
\]
\end{proof}

\section*{Main result}

\begin{theorem}
Let \(K \subset S^3\) be an L--space knot.  For every slope
\(r \in \mathbb Q \cup \{\infty\}\), the manifold \(S^3_r(K)\) has at most two
prime summands.
\end{theorem}

\begin{proof}
If \(K\) is the unknot, then every nonzero surgery is a lens space, \(0\)-surgery
is \(S^1 \times S^2\), and \(\infty\)-surgery is \(S^3\).  Each has at most one
prime summand.

Assume now that \(K\) is nontrivial.  Let
\[
M = S^3_r(K).
\]
If \(M\) is prime, there is nothing to prove.  Suppose therefore that \(M\) is
not prime.  Then \(M\) is reducible, so by Theorem~\ref{thm:greene}, \(K\) is a
\((p,q)\)-cable \(J_{p,q}\) of some knot \(J\), with \(p>1\), and the surgery
slope is the cabling slope \(r=pq\).  Lemma~\ref{lem:cable} gives
\[
M \cong L(p,q) \# S^3_{q/p}(J).
\]

The first factor \(L(p,q)\) is prime by Lemma~\ref{lem:lens-prime}.  Since
\(\gcd(p,q)=1\) and \(p>1\), the slope \(q/p\) is nonintegral.  Hence
\(S^3_{q/p}(J)\) is irreducible by Lemma~\ref{lem:nonintegral}, except in the
inessential case where it is \(S^3\).  Therefore the displayed connected-sum
decomposition has at most two non-\(S^3\) prime factors.  By the
Kneser--Milnor uniqueness theorem for prime decompositions \cite{Milnor},
\(M\) has at most two prime summands.
\end{proof}

\begin{thebibliography}{9}

\bibitem{GL}
C.~McA. Gordon and J.~Luecke,
\emph{Reducible manifolds and Dehn surgery},
Topology \textbf{35} (1996), no.~2, 385--409.

\bibitem{Gor83}
C.~McA. Gordon,
\emph{Dehn surgery and satellite knots},
Trans. Amer. Math. Soc. \textbf{275} (1983), no.~2, 687--708.

\bibitem{Greene}
J.~E. Greene,
\emph{L-space surgeries, genus bounds, and the cabling conjecture},
J. Differential Geom. \textbf{100} (2015), no.~3, 491--506.

\bibitem{Hempel}
J.~Hempel,
\emph{3-Manifolds},
Annals of Mathematics Studies, No.~86,
Princeton University Press, Princeton, NJ, 1976.

\bibitem{Milnor}
J.~Milnor,
\emph{A unique decomposition theorem for 3-manifolds},
Amer. J. Math. \textbf{84} (1962), 1--7.

\end{thebibliography}

\end{document}